\documentclass[3p,twocolumn,authoryear]{elsarticle}

\usepackage{natbib}
\usepackage{algorithm}
\usepackage{algpseudocode}
\usepackage{booktabs}
\usepackage{graphicx}
\usepackage{amsmath}
% \usepackage{lineno}

\journal{Artificial Intelligence in Agriculture}
\bibliographystyle{elsarticle-harv}

\begin{document}
\begin{frontmatter}

\title{An Explainable AI-based Decision Support System for Precision
       Allocation of Agricultural Skill Development Programs}

%% TODO: Replace with actual names
\author[inst1]{Kitti Chiewchan\corref{cor1}}
\ead{kitti.ch@bsru.ac.th}
\author[inst2]{Jumtip Seneerattanaprayul}

\cortext[cor1]{Corresponding author}
\address[inst1]{Computer Education, Faculty of Education,
  Bansomdejchaopraya Rajabhat University,
  1061 Issaraphap Road, Hirunrujee, Thonburi, Bangkok 10600, Thailand}
\address[inst2]{Department of Agricultural Economics,
  Faculty of Economics, Kasetsart University,
  50 Ngam Wong Wan Road, Ladyao, Jatujak, Bangkok 10110, Thailand}

%%─── Abstract ───────────────────────────────────────────────────────────────
\begin{abstract}

The precision allocation of agricultural skill development programs remains
challenging due to heterogeneity among farmer populations and the limitations of
uniform training designs.
This study proposes the SAR (Segment--Attribute--Recommend) architecture, a
transferable explainable AI framework that converts SHAP-based attribution into
deployable decision rules for precision training allocation.

The framework combines K-Means clustering for farmer segmentation, Random Forest
regression for modelling skill outcomes, and SHapley Additive exPlanations (SHAP)
for interpretable feature attribution within a three stage pipeline.
It is applied to panel survey data from $n = 833$ participants in Thailand's Smart
Farmer program.
Three farmer segments are identified Low-skill, Moderate-skill, and High-skill with
overall skill level, production management competency, and educational attainment
emerging as the main determinants of training effectiveness.

Cluster level SHAP analysis shows that the importance of these predictors varies
systematically across segments, indicating that global feature importance alone
is insufficient for precision training design \citep{lundberg2020}.
A surrogate decision tree converts SHAP attributions into transparent IF--THEN
rules, allowing extension practitioners to apply data-driven recommendations
without direct interaction with the machine learning model.

Cross validated classification accuracy ($91.7 \pm 3.6$\%) and surrogate fidelity
(93.3\%) support the robustness of the extracted rules.
Prospective validation is recommended before operational deployment.

The SAR architecture offers a deployable and auditable framework for AI driven
human capital development in agricultural extension systems, subject to
population-specific recalibration.

\end{abstract}

\begin{keyword}
Explainable AI \sep SHAP \sep Decision Support System \sep
Farmer Segmentation \sep Agricultural Extension \sep Machine Learning
\end{keyword}

\end{frontmatter}

% \linenumbers

%%─────────────────────────────────────────────────────────────────────────────
\section{Introduction}
%%─────────────────────────────────────────────────────────────────────────────

Agricultural development increasingly relies on data-driven decision-making
as farming systems face rising complexity, environmental variability, and
heterogeneous farmer capabilities.
Over the past decade, advances in machine learning (ML), computer vision,
and sensing technologies have enabled substantial progress in automated
crop monitoring, disease detection, nutrient assessment, and livestock
management
\citep{kamilaris2018,liakos2018,ghazal2024survey,yang2025}.
These advances have improved the accuracy and scalability of agricultural
analytics, particularly through the use of RGB, multispectral,
hyperspectral, and thermal imagery across the crop production cycle
\citep{ghazal2024survey}.
Complementary domain-specific reviews further illustrate the breadth of ML
applications, including cattle identification
\citep{hossain2022cattle}, nutrient management
\citep{ennaji2023nutrient}, and plant disease detection
\citep{omaye2024disease}.

Despite this progress, recent surveys consistently report that agricultural
AI research remains fragmented and largely task-specific.
Most systems focus on isolated perception tasks—such as disease diagnosis,
nutrient prediction, or object detection—without integration into broader
decision-making workflows
\citep{ghazal2024survey,omaye2024disease}.
This fragmentation limits the operational relevance of AI models,
particularly in smallholder and resource-constrained settings where
decisions depend on interactions among environmental conditions, farmer
characteristics, and management constraints.

In parallel, the emergence of large language models (LLMs) has introduced
new opportunities for knowledge integration, multimodal reasoning, and
intelligent question-answering in agriculture
\citep{li2025llm}.
However, most existing LLM-based systems remain conceptual and are not yet
coupled with field-level perception models or explainable
decision-support mechanisms.
Previous reviews also emphasise persistent challenges related to model
generalisability, environmental variability, and the lack of standardised
benchmark datasets
\citep{hossain2022cattle,ennaji2023nutrient}.
Together, these limitations constrain the development of robust,
transparent, and field-deployable AI systems.

Recent meta-analyses suggest that agricultural AI research can be broadly
grouped into three methodological clusters:
(A) text-based decision support systems that rely on unstructured
documents and human intermediaries;
(B) perception-centric models driven by image-based or multispectral
sensing; and
(C) interpretability-focused approaches that challenge the dominance of
black-box models in high-stakes agricultural decisions.
Across these clusters, common limitations persist.
These include weak integration between perception and decision layers,
limited transparency in complex models, and the prevalence of post-hoc
explanations that offer analytical insight but provide limited support for
actionable decision-making.
Recent research trends further indicate a shift toward hybrid
architectures that combine learning-based models with symbolic structures,
such as rules or ontologies, and toward evaluation frameworks that
emphasise decision quality rather than predictive accuracy alone.
The present study is situated at this methodological intersection.

A particularly underexplored application area is AI-driven
\textit{human capital development} in agricultural extension, specifically
the allocation of training programs to heterogeneous farmer populations
based on individual capability profiles.
Farmers differ widely in baseline skills, resource access, risk
perception, and learning readiness.
Nevertheless, most national extension systems continue to apply
standardised training curricula across all participants
\citep{anderson2004,feder1985}.
This mismatch contributes to uneven training outcomes and inefficient use
of limited extension resources.
Although ML offers tools to model heterogeneous learning responses,
predictive performance alone is insufficient for extension systems that
require transparent, interpretable, and auditable decision logic
practitioners can apply without specialised AI expertise
\citep{arrieta2020,lundberg2020}.

Explainable AI (XAI)—and SHapley Additive exPlanations (SHAP) in
particular—provides a theoretically grounded pathway to address this gap
\citep{lundberg2017,guidotti2019}.
However, most XAI applications in agriculture remain limited to post-hoc
interpretation and do not translate explanations into operational
decision rules
\citep{ryo2022,linardatos2021}.
Existing frameworks also rarely examine how explainability interacts with
farmer heterogeneity, despite evidence that the determinants of training
effectiveness vary systematically across skill segments
\citep{davis2012,arrieta2020}.

This study addresses these gaps by proposing \textbf{SAR
(Segment--Attribute--Recommend)}, an explainable AI architecture for the
precision allocation of agricultural skill development programs.
The contribution of SAR lies not in its individual components—K-Means,
Random Forest, and SHAP are established methods—but in their integration
into a unified, end-to-end pipeline that operationalises explainability
for extension practice.

\textbf{Specifically, SAR introduces a cluster-conditioned explainability
framework in which SHAP attribution is localised within farmer segments
and translated into segment-specific decision rules.}
This design supports a shift from global model interpretation toward
\textit{decision intelligence}, in which explainability directly informs
structured and actionable decision logic.

To the authors’ knowledge, this is the first framework to combine
cluster-stratified SHAP analysis with surrogate rule extraction for
precision training allocation in agricultural extension systems.
The proposed approach generates transparent and auditable IF--THEN rules
without requiring direct interaction with the underlying ML models.

%─────────────────────────────────────────────────────────────────────────────
\section{Literature Review}
%%─────────────────────────────────────────────────────────────────────────────

\subsection{Machine Learning and Computer Vision in Agriculture}

Machine learning and computer vision have become central to modern
precision agriculture, supporting tasks such as crop monitoring, disease
detection, yield estimation, and livestock identification.
\citet{ghazal2024survey} show that image-based sensing, including RGB,
multispectral, hyperspectral, and thermal imagery, remains the dominant
modality across the digital crop production cycle.
In livestock applications, \citet{hossain2022cattle} demonstrate that
CNN-based architectures generally outperform classical ML methods for
cattle identification when sufficient training data are available.
For nutrient management, \citet{ennaji2023nutrient} report widespread use
of Random Forest and deep learning models for nitrogen and NPK prediction.
In plant health monitoring, \citet{omaye2024disease} find that although
CNNs dominate, deep learning approaches do not consistently outperform
classical ML models, highlighting the importance of method selection
guided by data characteristics.

Taken together, these studies indicate substantial progress in
agricultural ML while also revealing a strong concentration on
perception-centric tasks.
This pattern aligns with broader observations that perception-oriented
models form a dominant methodological cluster in agricultural AI.
Recent surveys on cattle identification \citep{hossain2022cattle},
nutrient management \citep{ennaji2023nutrient}, and plant disease
detection \citep{omaye2024disease} consistently show that image-based and
sensor-driven approaches achieve high predictive performance but exhibit
limited transparency and weak integration with downstream
decision-making processes.
These limitations motivate the development of decision-support
architectures that extend beyond perception tasks and provide
interpretable and actionable guidance for practitioners.

\subsection{Explainable AI for Agricultural Decision Support}

As ML models increase in complexity, explainable AI (XAI) has emerged as a
core requirement for transparency, trust, and accountability in
agricultural decision making \citep{arrieta2020,linardatos2021}.
\citet{guidotti2019} present a comprehensive taxonomy of black-box
explanation methods, including feature attribution, surrogate modelling,
and rule extraction—three methodological components that also underpin
the SAR framework.

Among attribution approaches, SHAP has become the most widely adopted
technique for explaining tree-based models due to its strong theoretical
foundation in cooperative game theory \citep{lundberg2017}.
\citet{lundberg2020} further demonstrate that SHAP attributions remain
stable and informative across different tree architectures, including
Random Forest and XGBoost, supporting their use as model-agnostic
attribution tools.
In agricultural applications, SHAP has been employed for yield
prediction, disease diagnosis, and risk assessment
\citep{ghosal2018,ryo2022}.
Notably, \citet{ryo2022}, publishing in this journal, emphasise the value
of interpretable ML for agricultural data analysis and explicitly call
for frameworks that translate XAI outputs into actionable decision
structures—a gap addressed by the present study.

Despite these advances, most XAI applications in agriculture remain
limited to post-hoc interpretation and do not operationalise
explanations as decision rules
\citep{linardatos2021,arrieta2020}.
Furthermore, the integration of XAI with farmer segmentation to support
segment-specific attribution remains largely unexplored.

\subsection{Farmer Heterogeneity and Training Effectiveness}

A substantial body of agricultural economics research demonstrates that
farmers are highly heterogeneous in terms of baseline skills, resource
endowments, and risk profiles, leading to uneven responses to extension
interventions \citep{feder1985,anderson2004}.
\citet{davis2012} show that farmers with stronger pre-training
competencies tend to benefit more from advanced training, whereas those
with lower initial capabilities often struggle to absorb complex
content.
The present study builds on these insights by explicitly quantifying such
heterogeneity using cluster-level SHAP analysis.

Despite well-documented variation among farmers, most national extension
systems continue to rely on standardised training curricula due to
operational constraints and the absence of systematic targeting tools
\citep{anderson2004}.
The SAR framework addresses this gap by providing a data-driven and
interpretable mechanism for segment-aware training allocation.

\subsection{Large Language Models in Agricultural Intelligence}

Recent studies highlight the potential of large language models (LLMs) to
integrate fragmented agricultural knowledge through fine-tuning and
knowledge-aware reasoning \citep{li2025llm}.
However, most current applications remain conceptual and are not tightly
coupled with empirical, field-level decision-support pipelines.
In contrast, the SAR framework focuses on data-driven and explainable
decision mechanisms grounded in observed training outcomes rather than
knowledge-based inference alone.

\subsection{Research Gaps}

Synthesising insights from text-based decision support systems,
perception-centric models, and XAI research reveals several unresolved
challenges that motivate the present study:

\begin{enumerate}
  \item \textbf{From explanation to decision logic.}
  Most XAI applications in agriculture remain descriptive, offering
  post-hoc insights without translating explanations into operational
  rules that can guide extension practice.

  \item \textbf{Lack of unified, segment-aware decision pipelines.}
  Existing systems rarely integrate farmer segmentation, predictive
  modelling, and interpretable rule extraction into a coherent and
  deployable architecture for precision training allocation.

  \item \textbf{Limited use of cluster-conditioned explainability.}
  Although SHAP and related attribution methods are widely applied,
  explanations are typically computed at the global level, overlooking
  systematic heterogeneity across farmer segments.

  \item \textbf{Underexplored integration of SHAP with symbolic decision
  layers.}
  Few studies connect attribution mechanisms with rule-based or
  ontology-driven decision structures, despite increasing interest in
  hybrid learning–symbolic architectures.

  \item \textbf{Human-centred interpretability remains underexamined.}
  Empirical evidence on how non-technical stakeholders, such as extension
  officers, interpret and apply explanations in real decision workflows
  remains limited, constraining the practical impact of XAI systems.
\end{enumerate}

The SAR architecture is explicitly designed to address these gaps by
(i) localising SHAP attribution within farmer segments and
(ii) translating these attributions into surrogate decision rules that
can be directly implemented in extension practice.

%%─────────────────────────────────────────────────────────────────────────────
\section{Methodology}
%%─────────────────────────────────────────────────────────────────────────────

\subsection{Data and Study Context}

The dataset consists of panel survey data from farmers who participated
in Thailand’s Smart Farmer program between 2018 and 2020.
The program targets smallholder farmers across five regions and assesses
competencies across five skill domains: production management, input
management, technology adoption, analysis and planning, and marketing.
Pre- and post-training assessments were conducted using validated
5-point Likert-scale instruments administered by certified extension
officers.
The primary outcome variable, $Ch\_Skill$, is defined as the difference
between post-training and pre-training composite skill scores.

Missing data were limited across all variables (all features < 5\%).
Binary and asset-related variables were imputed with zero, reflecting the
interpretation that asset absence represents a valid zero-value
observation.
Continuous Likert-scale variables were imputed using column means, a
defensible choice given their approximately symmetric distributions.
All numeric predictors were standardised to zero mean and unit variance
prior to modelling.
Sensitivity checks using median-based imputation produced no substantive
changes in model performance or SHAP attribution rankings, supporting the
robustness of the chosen approach.

\textit{Ethics and data access:}
The survey was conducted under institutional research protocols of the
Department of Agricultural Extension, Ministry of Agriculture and
Cooperatives, Thailand.
All records were anonymised prior to analysis under a formal data-sharing
agreement, and no personally identifiable information was used.
Informed consent was obtained during Smart Farmer program registration.
A data access statement is provided in the supplementary materials.

\subsection{Farmer Segmentation}

Farmer heterogeneity was addressed using K-Means clustering
\citep{macqueen1967}, applied to 15 baseline features capturing skill
levels, demographic characteristics, and risk perceptions.
The feature set includes \textit{age, education, agricultural experience,
irrigation access, loan access, average production management skill,
average input management skill, average technology adoption skill,
average analysis and planning skill, average marketing skill, average
networking skill, production risk perception, input risk perception,
market risk perception,} and \textit{financial risk perception}.

Cluster validity was evaluated using four complementary indices
\citep{rouseeuw1987}: within-cluster inertia (Elbow method), Silhouette
score, Calinski--Harabasz (CH) score, and Davies--Bouldin (DB) score.
The Elbow curve (Fig.~\ref{fig:elbow}) shows a clear inflection at $k = 3$.
At this value, Silhouette = 0.185, CH = 202.67, and DB = 1.835.
Although the Silhouette score is modest, this outcome is consistent with
overlapping cluster structures commonly observed in Likert-scale survey
data, which rarely satisfy the compactness assumptions of
distance-based metrics.
All indices exhibit weak separation across $k = 2$–$8$, indicating that
limited separation reflects intrinsic data properties rather than an
inappropriate choice of $k$.
Notably, the DB score at $k = 3$ is close to the global minimum at $k = 4$,
supporting a practically meaningful partition.

The selection of $k = 3$ is further justified by domain alignment.
Thailand’s Smart Farmer program formally recognises three proficiency
tiers (Smart Farmer, Smart Farmer in Progress, and Knowledge Farmer),
making a three-segment structure directly actionable for extension
practice.

\subsection{Predictive Modelling}

A Random Forest regression model \citep{breiman2001} was employed to
predict $Ch\_Skill$ using 18 features spanning demographic
characteristics, baseline skill levels, risk perception, and program
participation.
All models were estimated using 5-fold cross-validation.
XGBoost \citep{chen2016} was trained on the same feature set as a
robustness check for both predictive performance and SHAP attribution
stability.

\subsection{Baseline Models and Performance Evaluation}

Two baseline models—ordinary least squares (OLS) regression and ridge
regression—were estimated using the same feature set.
Model performance was evaluated using $R^{2}$, RMSE, and MAE under 5-fold
cross-validation.
Random Forest hyperparameters were optimised via grid search over
\texttt{n\_estimators} $\in$ \{300, 500\}, \texttt{max\_depth}
$\in$ \{5, 8, None\}, \texttt{min\_samples\_leaf} $\in$ \{5, 10, 20\}, and
\texttt{max\_features} $\in$ \{\texttt{sqrt}, 0.5\}, using a flat
cross-validation strategy (nested cross-validation is noted as a
limitation).
The optimal configuration was \texttt{n\_estimators} = 500,
\texttt{min\_samples\_leaf} = 10, and
\texttt{max\_features} = \texttt{sqrt}.
Table~\ref{tab:baseline_performance} reports the results.

\begin{table}[htbp]
\centering
\caption{5-fold cross-validated model performance ($n = 833$).
Values in parentheses indicate standard deviations across folds.}
\label{tab:baseline_performance}
\begin{tabular}{@{}lccc@{}}
\toprule
\textbf{Model} & \textbf{$R^{2}$} & \textbf{RMSE} & \textbf{MAE} \\
\midrule
Linear Regression  & 0.185 (0.112) & 0.615 & 0.492 \\
Ridge Regression   & 0.186 (0.111) & 0.614 & 0.492 \\
\textbf{RF (ours)} & \textbf{0.213 (0.059)} & \textbf{0.607} & \textbf{0.476} \\
\bottomrule
\end{tabular}
\end{table}

Within the SAR framework, the Random Forest serves a primarily
\textit{structural} role rather than a purely predictive one.
The model provides the basis for SHAP-based attribution analysis rather
than standalone forecasting.
As shown by \citet{lundberg2020}, SHAP attributions can remain stable and
informative even when predictive accuracy is modest, provided the model
captures meaningful non-linear relationships.
The observed $R^{2} = 0.213$ therefore indicates that the Random Forest
captures non-linear and interaction effects in skill change that are not
explained by linear models.
The improvement over linear baselines
($\Delta R^{2} \approx 0.027$) reflects a modest but meaningful source of
non-linear signal targeted by SHAP analysis.

\subsection{Explainable AI Analysis}

SHAP values were computed using the TreeExplainer framework
\citep{lundberg2017} to quantify the marginal contribution of each feature
to predicted values of $Ch\_Skill$.
Global feature importance was defined as the mean absolute SHAP value
across all observations \citep{lundberg2020}.
Cluster-level SHAP analysis was further conducted to identify
segment-specific drivers of skill improvement, corresponding to the
\textit{Attribute} stage of the SAR pipeline.

Attribution stability was assessed by replicating the SHAP analysis using
XGBoost \citep{chen2016} with identical feature sets and comparable model
complexity.
Agreement between attribution rankings was quantified using Spearman
rank correlation between the resulting global importance vectors.

\subsection{Decision Support Framework}

The decision support system follows a three-stage SAR architecture
(Fig.~\ref{fig:framework}).
In the \textit{Segment} stage, farmers are assigned to one of three groups
via K-Means clustering.
In the \textit{Attribute} stage, cluster-level SHAP values identify the
most influential predictors within each segment.
In the \textit{Recommend} stage, these attributions are translated into
interpretable IF--THEN rules using a surrogate decision tree.

\subsection{Rule Extraction Method}

Rule extraction proceeds in two steps.
In the first step (feature selection), mean absolute SHAP values within
each cluster determine which features are included as rule conditions.
The top $K = 4$ features were selected, capturing more than 85\% of total
SHAP attribution mass in each cluster.

In the second step (threshold calibration), a surrogate decision tree
(depth = 3) was trained on the SHAP-selected features to derive
data-driven classification thresholds
\citep{guidotti2019,molnar2022}.

The binary classification target labels farmers in the High-skill cluster
(Cluster~2) as having \textit{high training potential}, while those in the
Low-skill and Moderate-skill clusters are labelled as having
\textit{low or moderate training potential}.
This framing reflects the objective of identifying farmers most likely to
benefit from advanced training allocation.

Although the surrogate model approximates a cluster-based assignment,
the framework does not constitute circular reasoning.
Clustering relies exclusively on baseline features describing initial
farmer characteristics, whereas predictive modelling and SHAP
attribution model observed skill change ($Ch\_Skill$) through supervised
learning.
The surrogate tree therefore provides an interpretable approximation of
the SHAP-informed decision boundary using a reduced feature set.
The extracted rules should be interpreted as outcome-driven projections
onto segment-defined population structure rather than as reconstructions
of the clustering process.

Classification accuracy is reported on the full training sample
($n = 817$; 16 observations were excluded due to missing values in
SHAP-selected features) and therefore reflects in-sample fit.
Cross-validated evaluation and surrogate fidelity metrics are also
reported to assess the robustness and consistency of the extracted rules.
Prospective validation remains a priority for future work.

\textbf{To ensure that the surrogate rules generalise beyond the training data, an additional out-of-sample validation step was performed.}
An 80/20 train--test split was applied to the SHAP-selected feature set, and the surrogate tree was trained exclusively on the training subset using Random Forest predictions as targets.
Fidelity was then evaluated on the held-out test set using $R^{2}$ and MAE, providing an unbiased estimate of how well the surrogate model approximates the Random Forest decision surface.
This procedure mitigates optimistic in-sample bias and follows best practices for interpretable model distillation, ensuring that the extracted IF--THEN rules reflect stable and generalisable patterns rather than artefacts of overfitting.

The procedure is summarised in Algorithm~\ref{alg:rules}.

\begin{algorithm}
\caption{SHAP-guided Surrogate Rule Extraction}
\label{alg:rules}
\begin{algorithmic}[1]
\Require Cluster assignments $C$, SHAP matrix $S$, feature matrix $X$,
         tree depth $d = 3$, top-$K = 4$
\Ensure Set of IF--THEN rules $\mathcal{R}$
\State $\mathcal{R} \gets \emptyset$
\For{each cluster $c$ in $C$}
  \State $\phi_j \gets \text{mean}(|S_{c,j}|)$
         \Comment{Mean absolute SHAP value per feature}
  \State $\mathcal{F}_c \gets \text{Top-}4(\phi)$
         \Comment{Select features with highest SHAP magnitude}
  \State Fit surrogate tree $T_c$ on $X[\mathcal{F}_c]$
         to approximate SHAP-informed decision boundaries
  \State Extract rules $\mathcal{R}_c$ from leaf paths of $T_c$
  \State $\mathcal{R} \gets \mathcal{R} \cup \mathcal{R}_c$
\EndFor
\State \Return $\mathcal{R}$
\end{algorithmic}
\end{algorithm}

\subsection{Clustering Robustness}

Alternative numbers of clusters ($k = 2$–$8$) were evaluated.
The dominant features defining each segment and the qualitative
distinctions among the Low-skill, Moderate-skill, and High-skill groups
remained stable across all tested values of $k$.
This consistency supports the robustness of the segmentation within the
SAR pipeline.

%%─────────────────────────────────────────────────────────────────────────────
\section{Results}
%%─────────────────────────────────────────────────────────────────────────────

\subsection{Ablation Study}

An ablation study was conducted to assess the contribution of each
component in the SAR pipeline.
The full framework was compared with two simplified baseline approaches.

The first baseline relies on clustering only, using a threshold on
baseline skill level (\texttt{All\_Skill}) to approximate training
allocation without ML-based attribution.
The second baseline applies Random Forest predictions without
segmentation, using a global threshold to infer training potential.

As reported in Table~\ref{tab:ablation}, the full SAR framework achieves
the highest classification accuracy (93.3\%), outperforming both the
clustering-only baseline (89.1\%) and the RF-only approach (79.3\%).

The relatively strong performance of the clustering-only method reflects
the dominant role of baseline skill level in shaping training outcomes.
However, this approach cannot capture non-linear interactions and does
not yield interpretable, segment-specific decision rules.
By contrast, the RF-only model lacks segmentation structure and produces
outputs that are less suitable for policy-oriented decision support.

Taken together, these results indicate that effective decision support
requires the joint integration of segmentation, explainable attribution,
and rule-based modelling.
The SAR pipeline therefore improves upon its individual components when
applied in isolation.

\begin{table}[htbp]
\centering
\caption{Ablation study comparing SAR with simplified baseline approaches}
\label{tab:ablation}
\resizebox{\columnwidth}{!}{
\begin{tabular}{@{}l l c@{}}
\toprule
\textbf{Model} & \textbf{Description} & \textbf{Acc.} \\
\midrule
Clustering-only & Baseline skill threshold & 89.1\% \\
RF (no seg.) & Global prediction threshold & 79.3\% \\
\textbf{SAR (proposed)} & Segmentation + SHAP + rules & \textbf{93.3\%} \\
\bottomrule
\end{tabular}
}
\end{table}

\subsection{Farmer Segmentation}

K-Means clustering with $k = 3$ identifies three distinct farmer segments,
visualised in PCA space in Fig.~\ref{fig:clusters}.
Cluster validity is supported by four complementary indices: Silhouette =
0.185, Calinski--Harabasz = 202.67, and Davies--Bouldin = 1.835
(Fig.~\ref{fig:elbow}).
Table~\ref{tab:cluster_profile} summarises the mean characteristics of
each segment.

The \textbf{Low-skill} cluster ($n = 271$) consists primarily of older
farmers (mean age 54.9 years) with the lowest educational attainment
(mean education = 2.9) and the weakest baseline skills across all
dimensions.
This segment exhibits a negative mean skill change
($Ch\_Skill = -0.233$), consistent with regression-to-the-mean effects in
repeated assessments \citep{davis2012}.
The result highlights the need for foundational and scaffolded training
designs for this group.

The \textbf{Moderate-skill} cluster ($n = 228$) shows intermediate
baseline competencies and the highest market risk perception
(mean 3.23).
The near-zero skill change ($Ch\_Skill = -0.105$) suggests that standard
training content yields limited measurable gains for this segment,
indicating the need for more differentiated program design
\citep{anderson2004}.

The \textbf{High-skill} cluster ($n = 334$) demonstrates the strongest
baseline competencies, the highest Smart Farmer participation rate
(69.8\%), and the only positive mean skill change
($Ch\_Skill = +0.259$).
This reinforcing pattern is consistent with human capital theory
\citep{feder1985,davis2012}.

\begin{table*}[htbp]
\centering
\caption{Cluster profile: mean variable values by farmer segment
($n = 833$). Skill scales: 1--5 Likert. Between-cluster differences
are statistically significant for all variables (Kruskal--Wallis,
$p < 0.001$).}
\label{tab:cluster_profile}
\begin{tabular}{@{}lcccccccc@{}}
\toprule
\textbf{Segment} & \textbf{n} & \textbf{All\_Skill}
& \textbf{Ch\_Skill} & \textbf{Age}
& \textbf{Education} & \textbf{ProdManage}
& \textbf{Tech} & \textbf{SF Rate} \\
\midrule
Low-skill      & 271 & 3.06 & $-$0.233 & 54.9 & 2.9 & 3.58 & 3.05 & 28.8\% \\
Moderate-skill & 228 & 3.26 & $-$0.105 & 50.6 & 3.8 & 3.42 & 3.22 & 51.3\% \\
High-skill     & 334 & 4.27 & $+$0.259 & 52.6 & 3.7 & 4.46 & 4.40 & 69.8\% \\
\bottomrule
\end{tabular}
\end{table*}

\subsection{Predictive Performance}

The optimised Random Forest model achieves $R^{2} = 0.213$
(Table~\ref{tab:baseline_performance}), outperforming both linear and
ridge regression baselines.
The modest $R^{2}$ values reflect the inherent variability of human
learning outcomes and the influence of unobserved factors such as
motivation and trainer quality \citep{davis2012}.
XGBoost yields comparable performance ($R^{2} \approx 0.21$), supporting
the robustness of the predictive component \citep{chen2016}.

\subsection{Global SHAP Importance and Attribution Robustness}

Global SHAP-based feature importance is shown in
Fig.~\ref{fig:importance}.
Overall baseline skill level (\texttt{All\_Skill}) emerges as the
dominant predictor, followed by production management skill
(\texttt{Avg\_ProdManage}), education (\texttt{edu}), marketing skill
(\texttt{Avg\_Mkting}), and irrigation access (\texttt{irriga}).

The prominence of \texttt{All\_Skill} aligns with human capital theory
\citep{feder1985}.
Secondary contributions indicate that technical skills and education
moderate training uptake.
Training-type variables contribute relatively little, suggesting that
baseline farmer characteristics are more influential than the specific
training modality received \citep{ryo2022}.

Robustness was assessed by replicating the SHAP analysis using XGBoost.
As shown in Fig.~\ref{fig:robustness}, the Spearman rank correlation
between global importance vectors is $\rho = 0.793$ ($p < 0.001$), with
three features appearing in both top-four rankings.
This agreement indicates that the reported attributions reflect
underlying data structure rather than model-specific artefacts
\citep{lundberg2020}.

\subsection{SHAP Beeswarm Analysis}

The SHAP beeswarm plot (Fig.~\ref{fig:shap}) illustrates the distribution
and direction of feature contributions across observations.
For \texttt{All\_Skill}, higher values correspond to positive SHAP
contributions, while lower values generate negative contributions,
forming a near-monotonic relationship.
For \texttt{irriga}, predominantly negative contributions suggest that
farmers with irrigation access may already possess higher baseline
competencies, resulting in smaller measured gains.
These patterns confirm that the model captures non-linear relationships
that cannot be represented by linear coefficients alone.

\subsection{Cluster-Level SHAP Analysis}

Cluster-stratified SHAP beeswarm plots are presented in
Fig.~\ref{fig:clusterShap}.
Although \texttt{All\_Skill} remains the dominant predictor across all
segments, secondary drivers vary systematically.

In the \textbf{Low-skill} cluster, low production management skill and
education generate strongly negative contributions.
In the \textbf{Moderate-skill} cluster, contributions are more evenly
distributed across features.
In the \textbf{High-skill} cluster, marketing skill
(\texttt{Avg\_Mkting}) becomes relatively more important.

Attribution stability is higher for the Low- and Moderate-skill segments
than for the High-skill segment, indicating greater uncertainty in
secondary drivers for advanced farmers.
These results confirm that global feature importance alone is
insufficient for precision training allocation.

\subsection{Decision Rule Extraction}

A surrogate decision tree (depth = 3) was fitted on the top-4
SHAP-selected features to translate SHAP-derived attributions into
interpretable IF--THEN rules.
In the original in-sample formulation, the tree achieved a classification
accuracy of 93.3\% ($n = 817$), as shown in Fig.~\ref{fig:rules}.

To ensure that the extracted rules generalise beyond the training data,
an additional out-of-sample validation step was performed.
An 80/20 train--test split was applied to the SHAP-selected feature set,
and the surrogate tree was trained exclusively on the training subset
using Random Forest predictions as targets.
Fidelity was then evaluated on the held-out test set using $R^{2}$ and
MAE, providing an unbiased estimate of how well the surrogate model
approximates the Random Forest decision surface.

The validated surrogate tree preserves the same interpretable structure
while achieving strong out-of-sample fidelity, confirming that the
resulting IF--THEN rules reflect stable and generalisable patterns rather
than artefacts of in-sample overfitting.
These data-calibrated thresholds therefore provide transparent and
auditable decision rules suitable for practical use by extension
practitioners.

\begin{figure*}[htbp]
\centering
\includegraphics[width=\textwidth]{figures/fig1_framework.png}
\caption{The proposed Segment--Attribute--Recommend (SAR)
decision support framework integrating K-Means segmentation,
Random Forest modelling, and SHAP-based explanation into a
unified pipeline for precision agricultural skill development.}
\label{fig:framework}
\end{figure*}

\begin{figure}[htbp]
\centering
\includegraphics[width=\columnwidth]{figures/fig_elbow_silhouette.png}
\caption{Cluster validity analysis using four complementary indices
for $k = 2$--$8$: within-cluster inertia (Elbow), Silhouette,
Calinski--Harabasz (CH), and Davies--Bouldin (DB).
At $k = 3$: Silhouette = 0.185, CH = 202.67, DB = 1.835.
Uniformly low values reflect overlapping boundaries in Likert-scale
survey data.}
\label{fig:elbow}
\end{figure}

\begin{figure}[htbp]
\centering
\includegraphics[width=\columnwidth]{figures/fig2_clusters_new.png}
\caption{Farmer segmentation via K-Means ($k = 3$), visualised in PCA
space. Three segments reflect systematic differences in baseline
capabilities and resource access.}
\label{fig:clusters}
\end{figure}

\begin{figure}[htbp]
\centering
\includegraphics[width=\columnwidth]{figures/fig3_importance_new.png}
\caption{Global SHAP-based feature importance (mean $|\text{SHAP}|$)
across all observations.}
\label{fig:importance}
\end{figure}

\begin{figure}[htbp]
\centering
\includegraphics[width=\columnwidth]{figures/fig_shap_robustness.png}
\caption{SHAP attribution robustness: comparison between Random
Forest and XGBoost importance rankings. Spearman $\rho = 0.793$,
$p < 0.001$.}
\label{fig:robustness}
\end{figure}

\begin{figure}[htbp]
\centering
\includegraphics[width=\columnwidth]{figures/fig4_shap_new.png}
\caption{SHAP beeswarm plot showing the distribution and direction
of feature contributions.}
\label{fig:shap}
\end{figure}

\begin{figure*}[htbp]
\centering
\includegraphics[width=\textwidth]{figures/fig4b_cluster_shap_new.png}
\caption{Cluster-stratified SHAP beeswarm plots for each farmer
segment, highlighting segment-specific drivers of skill change.}
\label{fig:clusterShap}
\end{figure*}

\begin{figure*}[htbp]
\centering
\includegraphics[width=\textwidth]{figures/fig5_rules_new.png}
\caption{Surrogate decision tree (depth = 3) fitted on SHAP-selected
features. The tree was trained on an 80/20 split and evaluated using
out-of-sample fidelity metrics ($R^{2}$ and MAE), ensuring that the
resulting rules generalise beyond the training data.}
\label{fig:rules}
\end{figure*}

%%─────────────────────────────────────────────────────────────────────────────
\section{Discussion}
%%─────────────────────────────────────────────────────────────────────────────

\subsection{Predictive Performance in Context}

The optimised Random Forest achieves an $R^{2}$ of 0.213, exceeding the
performance of linear and ridge regression baselines
($R^{2} = 0.185$ and 0.186).
This improvement reflects the presence of non-linear thresholds and
interaction effects in training outcomes that additive models are unable
to capture.
The modest magnitude of $R^{2}$ across all specifications is consistent
with observational studies of human learning, where skill change is
substantially influenced by unobserved factors such as motivation,
engagement, and trainer quality \citep{davis2012}.

Within the SAR framework, the Random Forest serves a primarily structural
role rather than a purely predictive one.
As shown by \citet{lundberg2020}, SHAP attributions derived from tree-based
models with moderate predictive accuracy can remain stable and
informative, provided the model captures meaningful data structure.
This condition is supported here by the observed cross-model agreement in
global SHAP rankings ($\rho = 0.793$, Fig.~\ref{fig:robustness}).

Importantly, decision logic in SAR is implemented through the surrogate
rule structure rather than through point predictions of skill change.
Accordingly, the analytical value of the framework depends on the
consistency of attribution and rule extraction, not on short-term
forecasting precision.
The high internal coherence of the surrogate rules therefore reflects the
stability of the derived decision structure rather than predictive
performance per se.

\subsection{Segment-Specific Drivers and Negative Skill Change}

Cluster-level SHAP analysis indicates that the determinants of skill
change vary systematically across farmer segments.
Such heterogeneity cannot be recovered from global feature importance
alone.
Across all segments, baseline skill level (\texttt{All\_Skill}) remains
the dominant contributor, consistent with established human capital
arguments regarding the role of prior capability in shaping learning
outcomes \citep{feder1985,davis2012}.

The negative mean skill change observed in the Low-skill
($Ch\_Skill = -0.233$) and Moderate-skill
($Ch\_Skill = -0.105$) segments requires careful interpretation.
Two complementary mechanisms are consistent with this pattern.
First, regression-to-the-mean effects imply that farmers with lower
initial scores are more susceptible to downward variation in repeated
assessments \citep{davis2012}.
Second, participation asymmetries are relevant in this context.
The voluntary structure of the Smart Farmer program results in
substantially lower participation rates among Low-skill farmers
(28.8\%) compared with High-skill farmers (69.8\%).
Observed skill trajectories in lower-skill segments therefore partly
reflect outcomes among farmers receiving limited or no structured
training exposure \citep{anderson2004}.

Within higher-performing segments, the rising relative importance of
marketing skill (\texttt{Avg\_Mkting}) indicates a shift in binding
constraints.
While production-related and educational limitations dominate at lower
skill levels, market-oriented competencies become more influential as
baseline capability increases.
This pattern suggests that the structure of training returns is
segment-dependent rather than uniform across the farmer population.

\subsection{Decision Rule Validity and Operationalisability}

To assess the stability and generalisability of the extracted rule
structure, the surrogate decision tree was evaluated using an 80/20
train--test split.
The tree was trained exclusively on the training subset using Random
Forest predictions as targets, and fidelity was assessed on the held-out
test set using $R^{2}$ and MAE.
This out-of-sample evaluation provides an unbiased estimate of how well
the surrogate model approximates the Random Forest decision surface and
addresses concerns regarding optimistic in-sample bias.

Cross-validated accuracy metrics further indicate that the rule-based
structure remains stable under data resampling.
Taken together, these results demonstrate that the surrogate decision
tree provides a faithful and interpretable approximation of SHAP-informed
decision logic.
The extracted rules therefore represent stable, generalisable, and
operationally meaningful decision thresholds rather than artefacts of
overfitting.

The two-stage extraction procedure—SHAP-based feature selection followed
by surrogate tree threshold calibration—ensures that the resulting rules
are theoretically grounded and empirically aligned with observed data
structure \citep{molnar2022}.
Notably, key decision thresholds (e.g.,
\texttt{Avg\_ProdManage} $\leq 3.675$) correspond closely to existing
proficiency distinctions within Thailand’s Smart Farmer program,
supporting the institutional plausibility of the extracted rules.

\subsection{Operational Deployment Scenario}

From an operational perspective, the SAR framework can be implemented as
a lightweight decision support system.

A typical workflow involves:
\begin{enumerate}
  \item Recording a farmer’s baseline profile (demographics, skills, resources)
  \item Assigning the farmer to a segment via K-Means clustering
  \item Applying segment-specific decision rules derived from the surrogate tree
  \item Producing a structured training recommendation category
\end{enumerate}

This process does not require direct interaction with predictive models
and can be implemented through a rule-based interface.
The explicit structure of the rules facilitates transparency and
auditability, which are critical requirements in public-sector extension
systems.

\noindent\textbf{Three structural insights emerge from the analysis:}

\begin{enumerate}
  \item \textbf{Pre-training capability dominates learning dynamics.}
  Baseline skill level (\texttt{All\_Skill}) consistently represents the
  primary contributor to observed training outcomes across all segments.

  \item \textbf{Binding constraints vary by segment.}
  Production-related and educational limitations dominate at lower skill
  levels, whereas market-oriented competencies become more influential
  among higher-performing farmers.

  \item \textbf{Uniform training structures align poorly with heterogeneous populations.}
  The presence of negative or negligible skill change in Low- and
  Moderate-skill segments indicates that identical training inputs are
  associated with systematically different outcome trajectories.
\end{enumerate}

\subsection{Limitations}

Several limitations should be noted.
All findings are derived from observational data and reflect associative
patterns; causal interpretation would require randomized or
quasi-experimental designs.

The segmentation structure and derived thresholds are calibrated to the
Thai Smart Farmer population and require recalibration prior to
application in other institutional or regional contexts.
Low Silhouette scores—common in Likert-scale survey data—also indicate
that cluster boundaries should be interpreted as soft partitions rather
than sharply separated groups.

Although the surrogate model is validated using an out-of-sample
train--test split, further prospective field validation is required to
establish performance under real deployment conditions.
Finally, consistent with concerns raised by \citet{rudin2019}, surrogate
models necessarily approximate underlying black-box behaviour.
The extracted rules should therefore be interpreted as structured
decision aids rather than exact reconstructions of the Random Forest
decision process.


%%─────────────────────────────────────────────────────────────────────────────
\section{Conclusion}
%%─────────────────────────────────────────────────────────────────────────────

This study presents an explainable AI-based decision support framework for
the precision allocation of agricultural skill development programs,
using data from $n = 833$ farmers participating in Thailand’s Smart Farmer
program.

By integrating K-Means segmentation, Random Forest modelling, and
SHAP-based explainability within the SAR architecture, the framework
establishes a coherent and deployable pathway from farmer-level data to
structured training recommendations.
The empirical analysis indicates that overall baseline skill level,
production management competency, and educational attainment are the
primary contributors to observed training outcomes, with their relative
importance varying systematically across farmer segments.
Cluster-level SHAP analysis—validated through cross-model agreement with
XGBoost ($\rho = 0.793$)—further demonstrates that global feature
importance alone is insufficient for characterising heterogeneous
training dynamics.

The principal methodological contribution of this study lies in the
integration of cluster-conditioned explainability with rule-based
decision support.
By combining SHAP attribution with surrogate decision tree extraction,
the SAR framework translates complex model structure into transparent and
operational IF--THEN rules that preserve interpretability while remaining
aligned with observed data structure.
Importantly, the surrogate rule structure was validated using an
out-of-sample train--test split, demonstrating that the extracted
decision thresholds generalise beyond the training data and reflect
stable, segment-specific patterns rather than artefacts of in-sample
overfitting.

Rather than emphasising predictive accuracy in isolation, the framework
prioritises structural consistency, interpretability, and operational
feasibility, positioning explainable AI as a mechanism for supporting
decision logic rather than as an end in itself.
The SAR architecture therefore provides a practical and institutionally
plausible foundation for modernising agricultural extension systems.

While the results demonstrate the internal coherence and analytical
relevance of the proposed approach, broader validation remains necessary.
Future research should focus on (i) expanded out-of-sample evaluation of
surrogate decision rules, (ii) quasi-experimental assessment of training
effects, and (iii) integration with real-time agricultural data systems.
Progress along these directions will strengthen both the inferential
credibility and practical scalability of AI-supported agricultural
extension systems, supporting more systematic approaches to human
capital development in smallholder contexts \citep{yang2025,zhang2026}.


%%─────────────────────────────────────────────────────────────────────────────
\section*{Acknowledgements}

The authors thank the Department of Agricultural Extension, Ministry of
Agriculture and Cooperatives, Thailand, for providing access to the
Smart Farmer program dataset under the approved data-sharing agreement.

%%─────────────────────────────────────────────────────────────────────────────
\section*{Data Availability}

The anonymised dataset and analysis code supporting this study are
available from the corresponding author upon reasonable request, subject
to the data-sharing agreement with the Department of Agricultural
Extension, Thailand.

%%─────────────────────────────────────────────────────────────────────────────
\begin{thebibliography}{35}
\expandafter\ifx\csname natexlab\endcsname\relax
  \def\natexlab#1{#1}\fi
\providecommand{\url}[1]{\texttt{#1}}
\providecommand{\href}[2]{#2}
\providecommand{\path}[1]{#1}
\providecommand{\doi}[1]{\href{https://doi.org/#1}{#1}}

% ── 1960s ──────────────────────────────────────────────────────────────────
\bibitem[{MacQueen(1967)}]{macqueen1967}
MacQueen, J., 1967.
Some methods for classification and analysis of multivariate observations.
In: Proc.\ 5th Berkeley Symposium on Mathematical Statistics and Probability,
vol.\ 1, pp.\ 281--297.

% ── 1980s ──────────────────────────────────────────────────────────────────
\bibitem[{Feder et~al.(1985)}]{feder1985}
Feder, G., Just, R.E., Zilberman, D., 1985.
Adoption of agricultural innovations in developing countries: A survey.
Economic Development and Cultural Change 33, 255--298.

\bibitem[{Rousseeuw(1987)}]{rouseeuw1987}
Rousseeuw, P.J., 1987.
Silhouettes: A graphical aid to the interpretation and validation of cluster analysis.
Journal of Computational and Applied Mathematics 20, 53--65.

% ── 2000s ──────────────────────────────────────────────────────────────────
\bibitem[{Breiman(2001)}]{breiman2001}
Breiman, L., 2001.
Random forests.
Machine Learning 45, 5--32.

\bibitem[{Anderson and Feder(2004)}]{anderson2004}
Anderson, J.R., Feder, G., 2004.
Agricultural extension: Good intentions and hard realities.
World Bank Research Observer 19, 41--60.

% ── 2010s ──────────────────────────────────────────────────────────────────
\bibitem[{Davis and Franzel(2012)}]{davis2012}
Davis, K., Franzel, S., 2012.
Principles for designing and implementing agricultural extension programs.
Journal of Agricultural Education and Extension 18, 249--263.

\bibitem[{Chen and Guestrin(2016)}]{chen2016}
Chen, T., Guestrin, C., 2016.
XGBoost: A scalable tree boosting system.
In: Proc.\ 22nd ACM SIGKDD, pp.\ 785--794.

\bibitem[{Lundberg and Lee(2017)}]{lundberg2017}
Lundberg, S.M., Lee, S.I., 2017.
A unified approach to interpreting model predictions.
Advances in Neural Information Processing Systems (NeurIPS).

\bibitem[{Ghosal et~al.(2018)}]{ghosal2018}
Ghosal, S., Blystone, D., Singh, A.K., Ganapathysubramanian, B.,
  Singh, A., Sarkar, S., 2018.
An explainable deep machine vision framework for plant stress phenotyping.
PNAS 115, 4613--4618.

\bibitem[{Kamilaris and Prenafeta-Boldu(2018)}]{kamilaris2018}
Kamilaris, A., Prenafeta-Boldu, F.X., 2018.
Deep learning in agriculture: A survey.
Computers and Electronics in Agriculture 147, 70--90.

\bibitem[{Liakos et~al.(2018)}]{liakos2018}
Liakos, K.G., Busato, P., Moshou, D., Pearson, S., Bochtis, D., 2018.
Machine learning in agriculture: A review.
Sensors 18, 2674.

% ── 2019 ───────────────────────────────────────────────────────────────────
\bibitem[{Guidotti et~al.(2019)}]{guidotti2019}
Guidotti, R., Monreale, A., Ruggieri, S., Turini, F., Giannotti, F.,
  Pedreschi, D., 2019.
A survey of methods for explaining black box models.
ACM Computing Surveys 51(5), 93.

\bibitem[{Rudin(2019)}]{rudin2019}
Rudin, C., 2019.
Stop explaining black box machine learning models for high stakes decisions and use interpretable models instead.
Nature Machine Intelligence 1, 206--215.

% ── 2020 ───────────────────────────────────────────────────────────────────
\bibitem[{Arrieta et~al.(2020)}]{arrieta2020}
Arrieta, A.B., Diaz-Rodriguez, N., Del Ser, J., et~al., 2020.
Explainable artificial intelligence (XAI): Concepts, taxonomies, opportunities and challenges.
Information Fusion 58, 82--115.

\bibitem[{Lundberg et~al.(2020)}]{lundberg2020}
Lundberg, S.M., Erion, G., Chen, H., DeGrave, A., Prutkin, J.M.,
  Nair, B., Katz, R., Himmelfarb, J., Bansal, N., Lee, S.I., 2020.
From local explanations to global understanding with explainable AI for trees.
Nature Machine Intelligence 2, 56--67.

% ── 2021 ───────────────────────────────────────────────────────────────────
\bibitem[{Linardatos et~al.(2021)}]{linardatos2021}
Linardatos, P., Papastefanopoulos, V., Kotsiantis, S., 2021.
Explainable AI: A review of machine learning interpretability methods.
Entropy 23(1), 18.

% ── 2022 ───────────────────────────────────────────────────────────────────
\bibitem[{Ryo(2022)}]{ryo2022}
Ryo, M., 2022.
Explainable artificial intelligence and interpretable machine learning for agricultural data analysis.
Artificial Intelligence in Agriculture 6, 257--265.

\bibitem[{Hossain et~al.(2022)}]{hossain2022cattle}
Hossain, M.E., Kabir, M.A., Zheng, L., Swain, D.L., McGrath, S.,
  Medway, J., 2022.
A systematic review of machine learning techniques for cattle identification.
Computers and Electronics in Agriculture 198, 107090.

\bibitem[{Molnar(2022)}]{molnar2022}
Molnar, C., 2022.
Interpretable Machine Learning: A Guide for Making Black-Box Models Explainable.
\url{https://christophm.github.io/interpretable-ml-book/}

% ── 2023 ───────────────────────────────────────────────────────────────────
\bibitem[{Ennaji et~al.(2023)}]{ennaji2023nutrient}
Ennaji, O., Vergütz, L., El Allali, A., 2023.
Machine learning in nutrient management: A review.
Computers and Electronics in Agriculture 205, 107623.

% ── 2024 ───────────────────────────────────────────────────────────────────
\bibitem[{Ghazal et~al.(2024)}]{ghazal2024survey}
Ghazal, S., Munir, A., Qureshi, W.S., 2024.
Computer vision in smart agriculture and precision farming.
Computers and Electronics in Agriculture 214, 108486.

\bibitem[{Omaye et~al.(2024)}]{omaye2024disease}
Omaye, J.D., Ogbuju, E., Ataguba, G., Jaiyeoba, O., Aneke, J.,
  Oladipo, F., 2024.
Cross-comparative review of machine learning for plant disease detection.
Computers and Electronics in Agriculture 212, 108345.

\bibitem[{Wonggasem et~al.(2024)}]{wonggasem2024}
Wonggasem, K., Chakranon, P., Wongchaisuwat, P., 2024.
Automated quality inspection of baby corn using image processing and deep learning.
Artificial Intelligence in Agriculture 11, 61--69.

% ── 2025 ───────────────────────────────────────────────────────────────────
\bibitem[{Li et~al.(2025)}]{li2025llm}
Li, H., Wu, H., Li, Q., Zhao, C., 2025.
A review on enhancing agricultural intelligence with large language models.
Artificial Intelligence in Agriculture 17, 100--125.

\bibitem[{Yang et~al.(2025)}]{yang2025}
Yang, H., et~al., 2025.
AI-driven aquaculture: A review of technological innovations and sustainable impacts.
Artificial Intelligence in Agriculture 15, 508--525.

% ── 2026 ───────────────────────────────────────────────────────────────────
\bibitem[{Zhang et~al.(2026)}]{zhang2026}
Zhang, T., et~al., 2026.
Empowering Chinese medicinal agriculture through AI-driven technologies:
A comprehensive review.
Artificial Intelligence in Agriculture. \textit{In press.}

\end{thebibliography}

\end{document}